\documentclass[10pt,a4paper]{amsart}
\usepackage{amsmath,amssymb,amsthm,graphicx}
\usepackage{verbatim}
\newenvironment{reference}{\comment}{\endcomment}
%\newenvironment{reference}{}{}
\newenvironment{slogan}{\comment}{\endcomment}
\newenvironment{history}{\comment}{\endcomment}

% For commutative diagrams we use Xy-pic
\usepackage[all]{xy}

% We use 2cell for 2-commutative diagrams.
\xyoption{2cell}
% Theorem environments.
%
\theoremstyle{plain}
\newtheorem{theorem}[subsection]{Theorem}
\newtheorem{proposition}[subsection]{Proposition}
\newtheorem{lemma}[subsection]{Lemma}

\theoremstyle{definition}
\newtheorem{definition}[subsection]{Definition}
\newtheorem{example}[subsection]{Example}
\newtheorem{exercise}[subsection]{Exercise}
\newtheorem{situation}[subsection]{Situation}

\theoremstyle{remark}
\newtheorem{remark}[subsection]{Remark}
\newtheorem{remarks}[subsection]{Remarks}

\numberwithin{equation}{subsection}

% Macros
%
\def\lim{\mathop{\mathrm{lim}}\nolimits}
\def\colim{\mathop{\mathrm{colim}}\nolimits}
\def\Spec{\mathop{\mathrm{Spec}}}
\def\Hom{\mathop{\mathrm{Hom}}\nolimits}
\def\Ext{\mathop{\mathrm{Ext}}\nolimits}
\def\SheafHom{\mathop{\mathcal{H}\!\mathit{om}}\nolimits}
\def\SheafExt{\mathop{\mathcal{E}\!\mathit{xt}}\nolimits}
\def\Sch{\mathit{Sch}}
\def\Mor{\mathop{\mathrm{Mor}}\nolimits}
\def\Ob{\mathop{\mathrm{Ob}}\nolimits}
\def\Sh{\mathop{\mathit{Sh}}\nolimits}
\def\NL{\mathop{N\!L}\nolimits}
\def\CH{\mathop{\mathrm{CH}}\nolimits}
\def\proetale{{pro\text{-}\acute{e}tale}}
\def\etale{{\acute{e}tale}}
\def\QCoh{\mathit{QCoh}}
\def\Ker{\mathop{\mathrm{Ker}}}
\def\Im{\mathop{\mathrm{Im}}}
\def\Coker{\mathop{\mathrm{Coker}}}
\def\Coim{\mathop{\mathrm{Coim}}}

% Boxtimes
%
\DeclareMathSymbol{\boxtimes}{\mathbin}{AMSa}{"02}

%
% Macros for moduli stacks/spaces
%
\def\QCohstack{\mathcal{QC}\!\mathit{oh}}
\def\Cohstack{\mathcal{C}\!\mathit{oh}}
\def\Spacesstack{\mathcal{S}\!\mathit{paces}}
\def\Quotfunctor{\mathrm{Quot}}
\def\Hilbfunctor{\mathrm{Hilb}}
\def\Curvesstack{\mathcal{C}\!\mathit{urves}}
\def\Polarizedstack{\mathcal{P}\!\mathit{olarized}}
\def\Complexesstack{\mathcal{C}\!\mathit{omplexes}}
% \Pic is the operator that assigns to X its picard group, usage \Pic(X)
% \Picardstack_{X/B} denotes the Picard stack of X over B
% \Picardfunctor_{X/B} denotes the Picard functor of X over B
\def\Pic{\mathop{\mathrm{Pic}}\nolimits}
\def\Picardstack{\mathcal{P}\!\mathit{ic}}
\def\Picardfunctor{\mathrm{Pic}}
\def\Deformationcategory{\mathcal{D}\!\mathit{ef}}
\usepackage{cleveref}
\crefname{originallemma}{original lemma}{original lemmas}
\Crefname{originallemma}{Original lemma}{Original lemmas}
\makeatletter
\newlabel{algebra-lemma-dominate}{{00IA}{source}{}{sourcepointer-1}{}}
\newlabel{algebra-lemma-dominate@cref}{{[originallemma][0][]00IA}{[1][0][]source}{}{}{}}
\makeatother
\crefname{originaldefinition}{original definition}{original definitions}
\Crefname{originaldefinition}{Original definition}{Original definitions}
\makeatletter
\newlabel{definition-valuative-criterion}{{03IX}{source}{}{sourcepointer-2}{}}
\newlabel{definition-valuative-criterion@cref}{{[originaldefinition][0][]03IX}{[1][0][]source}{}{}{}}
\makeatother
\crefname{originallemma}{original lemma}{original lemmas}
\Crefname{originallemma}{Original lemma}{Original lemmas}
\makeatletter
\newlabel{morphisms-lemma-generalizations-lift-flat}{{03HV}{source}{}{sourcepointer-3}{}}
\newlabel{morphisms-lemma-generalizations-lift-flat@cref}{{[originallemma][0][]03HV}{[1][0][]source}{}{}{}}
\makeatother
\title{The existence part of the valuative criterion may be tested with finite separable extensions}
\author{The Stacks Project Authors}
\date{stacks-project:spaces-morphisms.tex:local-snapshot-20261012; GNU FDL 1.2 or later}
\begin{document}
\maketitle
\setcounter{subsection}{1}
\noindent\textit{Editorial note.} All delivered mathematical text is the author's own, copied from the pinned Stacks Project chapter file spaces-morphisms.tex; no statement or proof step was rewritten, reordered or removed.\par\medskip
\noindent\textit{Editorial note.} The author states that only the implication (1) implies (2) has to be shown, and the proof carries that implication out in full; the converse direction is the formal restriction of the same condition.\par\medskip
\noindent\textit{Editorial note.} The three references of the counted spans are recorded as non-load-bearing navigation pointers; the argument that consumes them stays inside the counted proof span.\par\medskip
\noindent\textit{Editorial note.} Preamble: the three native excerpt spans of sources/stacks-spaces-morphisms/preamble.tex (lines 8-13, 14-18 and 154-217). The chapter file itself carries no preamble text.\par\medskip
\noindent\textit{Allowed standard prerequisites (editorial).}\begin{itemize}
\item The shape of the diagram in the valuative criterion. Definition definition-valuative-criterion, used at proof lines 8451-8458. Fixes the commutative square that has to be completed.
\item Generalizations lift along flat morphisms of schemes. Morphisms, Lemma morphisms-lemma-generalizations-lift-flat, used at proof lines 8465-8469. Produces the generalization p' of p mapping to the generic point of Spec(A').
\item Existence of a valuation ring dominating a given local subring. Algebra, Lemma algebra-lemma-dominate, used at proof lines 8501-8504. Supplies the valuation ring A\_1 inside kappa(q') that dominates B\_q.
\item Base change of surjective etale morphisms of algebraic spaces. used at proof lines 8459-8463. Makes Spec(A') x\_X U -> Spec(A') surjective etale, so that the point p exists.
\end{itemize}\medskip
\ifdefined\hypertarget\hypertarget{sourcepointer-1}{}\fi
\noindent\textit{Source reference.} Original lemma 00IA: focused-seeds/AGP/sources/stacks-snapshot-20261011/tags.txt, source lines 672--672. Navigation pointer into the Stacks Project; the standard fact it supplies is applied inside the counted proof span and the argument resting on it stays there.\par\medskip
\ifdefined\hypertarget\hypertarget{sourcepointer-2}{}\fi
\noindent\textit{Source reference.} Original definition 03IX: focused-seeds/AGP/sources/stacks-snapshot-20261011/tags.txt, source lines 4619--4619. Navigation pointer into the Stacks Project; the standard fact it supplies is applied inside the counted proof span and the argument resting on it stays there.\par\medskip
\ifdefined\hypertarget\hypertarget{sourcepointer-3}{}\fi
\noindent\textit{Source reference.} Original lemma 03HV: focused-seeds/AGP/sources/stacks-snapshot-20261011/tags.txt, source lines 4581--4581. Navigation pointer into the Stacks Project; the standard fact it supplies is applied inside the counted proof span and the argument resting on it stays there.\par\medskip
\begin{lemma}
\label{lemma-finite-separable-enough}
Let $S$ be a scheme.
Let $f : X \to Y$ be a morphism of algebraic spaces over $S$.
The following are equivalent
\begin{enumerate}
\item $f$ satisfies the existence part of the valuative criterion
as in Definition \ref{definition-valuative-criterion},
\item $f$ satisfies the existence part of the valuative criterion
as in Definition \ref{definition-valuative-criterion} modified by
requiring the extension $K'/K$ to be finite separable.
\end{enumerate}
\end{lemma}

\begin{proof}
We have to show that (1) implies (2). Suppose given a diagram
$$
\xymatrix{
\Spec(K') \ar[r] \ar[d] & \Spec(K) \ar[r] & X \ar[d] \\
\Spec(A') \ar[r] \ar[rru] & \Spec(A) \ar[r] & Y
}
$$
as in Definition \ref{definition-valuative-criterion} with $K \subset K'$
arbitrary. Choose a scheme $U$ and a surjective \'etale morphism $U \to X$.
Then
$$
\Spec(A') \times_X U \longrightarrow \Spec(A')
$$
is surjective \'etale. Let $p$ be a point of $\Spec(A') \times_X U$
mapping to the closed point of $\Spec(A')$. Let $p' \leadsto p$
be a generalization of $p$ mapping to the generic point of $\Spec(A')$.
Such a generalization exists because generalizations lift along flat
morphisms of schemes, see
Morphisms, Lemma \ref{morphisms-lemma-generalizations-lift-flat}.
Then $p'$ corresponds to a point of the scheme $\Spec(K') \times_X U$.
Note that
$$
\Spec(K') \times_X U
=
\Spec(K') \times_{\Spec(K)} (\Spec(K) \times_X U)
$$
Hence $p'$ maps to a point $q' \in \Spec(K) \times_X U$ whose
residue field is a finite separable extension of $K$. Finally,
$p' \leadsto p$ maps to a specialization $u' \leadsto u$ on the
scheme $U$. With all this notation we get the following diagram of
rings
$$
\xymatrix{
\kappa(p') & & \kappa(q') \ar[ll] & \kappa(u') \ar[l] \\
& \mathcal{O}_{\Spec(A') \times_X U, p} \ar[lu] & &
\mathcal{O}_{U, u} \ar[ll] \ar[u] \\
K' \ar[uu] & A' \ar[l] \ar[u] & A \ar[l] \ar'[u][uu]
}
$$
This means that the ring $B \subset \kappa(q')$ generated by
the images of $A$ and $\mathcal{O}_{U, u}$ maps to a subring
of $\kappa(p')$ contained in the image $B'$ of
$\mathcal{O}_{\Spec(A') \times_X U, p} \to \kappa(p')$.
Note that $B'$ is a local ring. Let $\mathfrak m \subset B$
be the maximal ideal. By construction $A \cap \mathfrak m$,
(resp.\ $\mathcal{O}_{U, u} \cap \mathfrak m$, resp.\ $A' \cap \mathfrak m$)
is the maximal ideal of $A$ (resp.\ $\mathcal{O}_{U, u}$, resp.\ $A'$).
Set $\mathfrak q = B \cap \mathfrak m$. This is a
prime ideal such that $A \cap \mathfrak q$ is the maximal ideal of $A$.
Hence $B_{\mathfrak q} \subset \kappa(q')$ is a local ring dominating
$A$. By
Algebra, Lemma \ref{algebra-lemma-dominate}
we can find a valuation ring $A_1 \subset \kappa(q')$
with field of fractions $\kappa(q')$
dominating $B_{\mathfrak q}$. The (local) ring map
$\mathcal{O}_{U, u} \to A_1$ gives a morphism
$\Spec(A_1) \to U \to X$
such that the diagram
$$
\xymatrix{
\Spec(\kappa(q')) \ar[r] \ar[d] & \Spec(K) \ar[r] & X \ar[d] \\
\Spec(A_1) \ar[r] \ar[rru] & \Spec(A) \ar[r] & Y
}
$$
is commutative. Since the fraction field of $A_1$ is $\kappa(q')$
and since $\kappa(q')/K$ is finite
separable by construction the lemma is proved.
\end{proof}

\end{document}
